\documentclass[11pt,a4paper]{article}

\usepackage[T1]{fontenc}
\usepackage[utf8]{inputenc}
\usepackage{lmodern}
\usepackage{amsmath}
\usepackage[margin=2.6cm]{geometry}
\usepackage{microtype}
\usepackage{booktabs}
\usepackage{enumitem}
\usepackage{xcolor}
\usepackage{listings}
\usepackage{titlesec}
\usepackage[hidelinks]{hyperref}

\definecolor{shade}{gray}{0.96}
\definecolor{rule}{gray}{0.45}

\lstset{
  basicstyle=\ttfamily\footnotesize,
  backgroundcolor=\color{shade},
  frame=none,
  breaklines=true,
  columns=fullflexible,
  xleftmargin=0.6em,
  framexleftmargin=0.6em,
  aboveskip=1em,
  belowskip=1em,
}

\titleformat{\section}{\large\bfseries}{\thesection.}{0.6em}{}
\titleformat{\subsection}{\normalsize\bfseries}{\thesubsection}{0.6em}{}

\setlist[itemize]{leftmargin=1.2em,itemsep=0.25em,topsep=0.4em}

\newcommand{\tx}[1]{\texttt{\footnotesize #1}}

\title{\textbf{pq402}\\[0.3em]
\large An API that charges an AI agent, checks it is allowed to buy,\\
and never learns which buyer it was}
\author{Manuel Guilherme Galmanus}
\date{5 August 2026 \\[0.4em]
\small Stellar Summit SP 2026 · Payments and Agent Tooling (SDF DevEx)\\
\small sub-lane 3A, Agentic Payments (x402 / MPP)\\[0.6em]
\small \url{https://github.com/Galmanus/pq402}}

\begin{document}
\maketitle

\begin{abstract}
\noindent
Autonomous software is starting to buy things: data, API calls, compute. The
protocols for that are arriving --- x402 revives HTTP's long-dormant
\texttt{402 Payment Required} and settles it on-chain --- but they answer only
one question, \emph{did the agent pay}. Metered services have a second,
\emph{is this agent allowed to buy this at all}, and today's answers to it are
an API key or an account, both of which stamp the buyer's identity onto every
request they will ever make.

This paper describes a working system that answers both, and answers the second
without learning who is asking. The eligibility check is a hash-based
zero-knowledge proof judged by a Soroban contract, not by the server; a
per-challenge nullifier is burned in contract storage, so single use is
enforced by consensus and survives the server being restarted, replaced, or
dishonest. In its unlinkable mode the proof publishes a commitment rather than
a credential identifier, with fresh blinding on every use: the server learns
that \emph{a member of the issuer's set} paid and nothing further, and two
payments by one credential share no public value, so they cannot be linked to
each other.

Everything reported here runs on Stellar testnet and every claim is backed by a
transaction hash. We also record four undocumented details of the x402 protocol
on Stellar that each cost a debugging round to find, and publish the fix as a
package so the next team does not pay for them again. The limits are stated
with the same precision as the results: the proof system is post-quantum, the
transaction carrying it is signed with Ed25519 and is not.
\end{abstract}

\vspace{0.5em}
\hrule
\vspace{1.2em}

\section{The gap}

An agent that buys something must be authorised to buy it. The two mechanisms
in use today both work by identifying the buyer.

An \textbf{API key} is a bearer identifier. Every call carries it, so the seller
accumulates a complete, joinable history of one buyer's behaviour --- which is
usually not what either party wanted, merely what was easy to build.

An \textbf{account} is worse for the buyer and better for the seller, and has the
same shape: eligibility is established by being known.

Neither is a failure of care. There has not been an alternative that a server
could check cheaply, that a client could produce quickly, and that did not
require the server to be trusted with the answer. What follows is an attempt at
one.

\subsection{Why now}

Three things had to be true at once, and only recently are:

\begin{itemize}
  \item \textbf{A payment rail an agent can drive.} x402 defines a 402 response
    carrying machine-readable payment requirements, and on Stellar the client
    signs Soroban \emph{authorization entries} rather than transaction
    envelopes --- so a facilitator assembles the transaction and pays the
    network fee. An agent needs the asset it intends to spend and no XLM at
    all.
  \item \textbf{A proof cheap enough to verify inside a smart contract.}
    Circle STARKs over the Mersenne-31 field are the cheapest hash-based proof
    system in production. Verification fits an ordinary Soroban contract call.
  \item \textbf{Somewhere to put the single-use guarantee that is not the
    server.} Contract storage. A nullifier burned there cannot be forgotten by
    restarting a process.
\end{itemize}

\section{What the system does}

A request to a paid route meets two independent gates. Neither is decided by
the server, and the server cannot overrule either.

\begin{center}
\begin{tabular}{@{}lll@{}}
\toprule
\textbf{gate} & \textbf{mechanism} & \textbf{who decides} \\
\midrule
payment & x402 \texttt{exact}, USDC over its Soroban Asset Contract & Stellar consensus \\
credential & hash-based STARK, nullifier burned in storage & Stellar consensus \\
\bottomrule
\end{tabular}
\end{center}

The flow, in the order it happens:

\begin{lstlisting}
agent -> GET /premium
      <- 402  payment-required header + credential challenge

agent    proves eligibility locally, 7 ms
agent -> POST /pq/unlock  { proof, publics }
server-> Soroban: verify the proof and burn the nullifier, one transaction
      <- pass, and the transaction hash of the burn

agent -> GET /premium  + PAYMENT-SIGNATURE + the pass
server-> facilitator: verify, then settle
      <- 200, the goods, and the settlement hash
\end{lstlisting}

Two orderings in that sequence are deliberate and worth naming.

\textbf{The credential is checked before the payment.} Charging someone and then
refusing them is worse than refusing them, and arranging that is the server's
fault, not the protocol's.

\textbf{Verification precedes settlement.} \texttt{settle} moves money, so a
payload that fails verification must never reach it. The package has a test
asserting that a rejected verification leaves \texttt{settle} uncalled, because
the ordering is the only thing standing between a bad payload and a real
transfer.

\section{The credential, and what the server learns}

\subsection{Identifying mode}

The relation proves, in zero knowledge, that the caller knows a secret
$s$ such that
\[
  \mathrm{leaf} = \mathrm{trunc}_8\big(\pi(s \,\|\, \mathrm{action})\big),
  \qquad
  \mathrm{nullifier} = \mathrm{trunc}_8\big(\pi(s \,\|\, \mathrm{round})\big)
\]
share that one secret, where $\pi$ is the Poseidon2 permutation over
Mersenne-31. The AIR constrains the secret cells of the two permutation blocks
to be equal in the same row; without that constraint a prover could present one
member's leaf beside another member's nullifier.

Both published values are \emph{truncated} to a digest. An earlier version
published the full sixteen-limb permutation state, and that was a break rather
than a wart: $\pi$ is a bijection and the context sits public beside it, so a
single inversion returned the secret. It was found by audit, demonstrated
against a real testnet proof in milliseconds, and fixed by publishing digests
--- which puts recovery back at $|F|^8 \approx 2^{248}$.

In this mode the leaf is a public input. The server therefore knows
\emph{which} credential paid, and two payments by one credential are trivially
linked. That is anonymous in the sense of not knowing your name; it is not
anonymous in the sense of not knowing which of the thousand you are.

\subsection{Unlinkable mode}

The second mode closes that. The relation publishes
\[
  C = \mathrm{compress}(\mathrm{leaf} \,\|\, \mathrm{blinder})
\]
with a fresh blinder drawn from system entropy on every use, and never the
leaf. Two proofs compose on that shared $C$:

\begin{itemize}
  \item the \textbf{acting} half proves that a valid credential acted under $C$,
    and burns its nullifier;
  \item the \textbf{membership} half proves that $C$ sits under the issuer's
    published Merkle root, through a private authentication path.
\end{itemize}

What the server learns is exactly: \emph{someone in the issuer's set paid}. Not
which member. And because the blinder is fresh, two payments by one credential
share no public value at all, so they cannot be linked to each other either.

Three orderings again, each for a reason:

\begin{itemize}
  \item the two commitments are compared \textbf{before} anything is spent ---
    two proofs are only about one credential if they share $C$, and checking
    that is free while the proofs cost two transactions;
  \item membership is verified \textbf{before} the acting half --- membership is
    a read-only call and acting burns a nullifier, so there is no reason to
    burn one for a caller who was never in the set;
  \item the claimed root is compared against the \textbf{issuer's} --- a
    membership proof is only ever a statement about the tree the prover chose,
    and without that check anyone can build a tree containing their own leaf.
\end{itemize}

\subsection{The challenge is not the server's to choose}

The acting contract derives the expected round from the ledger sequence:
\[
  \mathrm{epoch} = \left\lfloor \frac{\mathrm{ledger\_sequence}}{720} \right\rfloor
\]
and refuses any other. Freshness is therefore guaranteed by consensus rather
than asserted by whoever is asking, and one credential acts once per epoch ---
about an hour at Stellar's close time. A server cannot replay an old challenge
at a member, and a member cannot pick a convenient one.

\section{Four things about x402 on Stellar that are not written down}

Each of these cost a debugging round, and each produced an error message that
pointed somewhere other than the cause. They are encoded as behaviour in the
published package rather than as advice, because a comment does not stop anyone
from getting it wrong.

\begin{enumerate}
  \item \textbf{v2 carries the requirements in a \texttt{payment-required}
    header, not the body.} The body is read only when
    \texttt{x402Version == 1}. A server that publishes perfect requirements as
    JSON alone is invisible to a v2 client, which fails with
    \emph{Invalid payment required response} having never looked at them. Every
    402 a gated route emits needs the header --- including the one meaning
    ``you have a session but still owe payment'', which is the one everybody
    forgets.
  \item \textbf{The signed payload returns in \texttt{PAYMENT-SIGNATURE}},
    with \texttt{X-PAYMENT} the v1 name. Read both.
  \item \textbf{\texttt{extra.areFeesSponsored} must appear in the
    requirements} or the Stellar \texttt{exact} scheme refuses to build a
    payment at all. That flag is the facilitator's assertion, so it should be
    fetched from \texttt{/supported} at startup rather than claimed on the
    facilitator's behalf.
  \item \textbf{Stellar assets carry seven decimals, not six.} \texttt{0.10} is
    \texttt{1000000} base units. The six an EVM habit expects is a tenth of the
    intended price, settling quietly --- the worst kind of bug, because it
    works.
\end{enumerate}

A fifth, adjacent: the Stellar CLI's \texttt{--network} takes a configuration
\emph{name} (\texttt{testnet}) while x402 takes a CAIP-2 \emph{identifier}
(\texttt{stellar:testnet}). Feeding one to the other fails with
\emph{Invalid name}, four frames from anything that mentions networks.

\section{Spending limits that are rules, not promises}

A client-side cap --- \texttt{-{}-max 0.10} on the agent --- is a promise the
agent makes to itself. An agent that has been prompt-injected, misconfigured,
or handed a bad tool description simply stops making it, and nothing outside
the agent notices.

\texttt{agent-treasury} is the same limit as a rule. It is a Soroban contract
holding the agent's funds, enforcing two policies before any token moves:

\begin{itemize}
  \item a \textbf{contract allow-list}, so a stolen policy key cannot be pointed
    at a token the treasury was never meant to touch;
  \item a \textbf{rolling daily cap}, keyed off ledger time rather than a
    counter that calling more often can reset.
\end{itemize}

Both refusals happen in consensus. On testnet: a payment inside policy settles;
one over the cap is refused with error \texttt{\#3}; one naming a token outside
the allow-list is refused with \texttt{\#2}; and the remaining budget is
unchanged after both, because a refusal consumes nothing.

The budget is committed \emph{before} the transfer: if the transfer panics the
whole invocation reverts, so the ordering cannot leak budget, while committing
afterwards would leave a window in which a reentrant call still sees the old
total.

\section{Measurements}

All on Stellar testnet, and each number was measured rather than estimated.

\begin{center}
\begin{tabular}{@{}lr@{}}
\toprule
proving, identifying mode (median of 15 runs, process start included) & 7 ms \\
proving, range across those runs & 6--11 ms \\
proof size, identifying mode & 79{,}186 B \\
proof size, unlinkable mode (acting half) & 113{,}360 B \\
proof size, unlinkable mode (membership half) & 82{,}765 B \\
on-chain verification, unlinkable act (12 queries, blowup $2^7$) & 243{,}213{,}882 instr. \\
on-chain verification, identifying spend (40 queries, blowup $2$) & 265{,}797{,}092 instr. \\
Soroban per-transaction ceiling & 400{,}000{,}000 instr. \\
settlement fee, paid by the facilitator & 23{,}073 stroops \\
\bottomrule
\end{tabular}
\end{center}

The instruction rows are read from the submitted envelopes' declared Soroban
resources, which the CLI sets from simulation; the transactions are the ones
linked in the repository and both succeeded.

Two things fall out of them. The first is that the ceiling is real and close:
one verification consumes 61\% of the per-transaction budget, so the security
parameters are not a matter of taste. At roughly 20M instructions per query,
the cap alone admits something under twenty of them, and that is before the
membership half and the storage writes.

The second is a result we did not expect and would not have found without
measuring. The identifying gate asks 40 queries at blowup 2 and costs
\emph{more} --- 265.8M against 243.2M --- while being the weaker of the two by a
wide margin. Blowup buys soundness at a discount that queries do not: raising
the rate lengthens the committed codeword, which the prover pays for, while the
verifier's cost tracks the number of openings it has to check. On a chain that
meters the verifier and not the prover, few queries over a high blowup strictly
dominates. The legacy configuration had it backwards.

For scale, SDF's own privacy-pool prototype spends roughly 40M instructions per
BLS12-381 pairing. A Groth16 verification needs several, so a hash-based
verifier here costs on the order of twice a pairing-based one. That factor of
two is the whole price of dropping the trusted setup and the exposure to Shor.

The fee row is the one worth pausing on: the account charged the network fee is
the facilitator's, not the agent's. That is fee sponsorship working, and it is
why an agent wallet needs only the asset it intends to spend.

\subsection{Security parameters}

FRI soundness is approximately $\text{queries} \times \log_2(\text{blowup})$
bits before the proof-of-work grind, and it cannot exceed the field the
challenges are drawn from. Both halves of that sentence matter.

The identifying gate originally ran at 40 queries with blowup 2: about 48 bits
conjectured, and near 24 against a Grover-assisted adversary, since Fiat-Shamir
grinding falls to a square-root search. Twenty-four bits is a demonstration
parameter wearing a production label.

Raising it on that path was impossible rather than merely expensive: its
challenge field is a degree-3 extension of Mersenne-31, about 93 bits, so 46
quantum bits is that path's ceiling at any configuration. The gate now used
draws challenges from QM31, a degree-4 extension with a 124-bit ceiling, and
runs at 16 queries with blowup $2^6$: roughly 96 conjectured bits, near 48
quantum. The configuration is pinned inside the contract rather than passed by
the caller, because whoever picks the security level picks how hard their own
proof is to forge.

\section{Reusability}

The protocol work is published, not merely described.

\begin{lstlisting}
npm install x402-stellar-paywall
\end{lstlisting}

\begin{lstlisting}[language=]
const pay = await expressPaywall({ price: "0.10", payTo: RECIPIENT });
app.get("/weather", pay, (req, res) =>
  res.json({ temp: 24, settled_by: req.x402.transaction }));
\end{lstlisting}

That is the whole integration. Adding \texttt{credential: \{ ... \}} adds the
eligibility gate; adding \texttt{mode: "crowd"} to it makes that gate
unlinkable. Bindings exist for Express, Hono and FastAPI --- the three
frameworks the sub-lane names --- over a framework-free core with no runtime
dependencies.

It was proven somewhere other than its birthplace: \texttt{examples/second-app}
is a weather API with no credential logic in it at all, which declares the
package as an ordinary dependency, gates one route with one line, and was paid
by an agent on testnet.

\section{Related work}

None of the cryptography here is new, and it would be dishonest to present it
as such. Post-quantum anonymous credentials are a studied object with a
literature, and this system is an engineering result standing on it. The
comparison worth making is not against nothing; it is against the published
numbers.

Policharla, Westerbaan, Faz-Hern\'andez and Wood\footnote{Cryptology ePrint
2023/414, \texttt{ia.cr/2023/414}.} build post-quantum Privacy Pass from
post-quantum anonymous credentials, proving possession of a zkDilithium
signature --- a STARK-friendly variant of Dilithium2 --- inside a STARK. Their
tokens run 85--175\,KB at a 115-bit proof security level, generate in
0.3--5\,s, and verify in 20--30\,ms. That is the closest published relative of
what runs here, and it is the reference a reader should hold this against.

Chen et al.\footnote{ePrint 2024/650, PQCrypto 2023.} give the first
post-quantum Direct Anonymous Attestation from symmetric primitives alone, a
SPHINCS+-derived credential shown through an MPC-in-the-Head proof. Feneuil and
Rivain's CAPSS\footnote{ePrint 2025/061, with the SmallWood proof system at
ePrint 2025/1085.} builds signatures from arithmetization-oriented permutations
--- Rescue-Prime, Poseidon, Griffin, Anemoi --- and reports anonymous-credential
presentation proofs under 150\,KB at 128 bits. That is the same primitive family
this work uses, and the better-optimised end of it. On the lattice side, Argo et
al.\footnote{ePrint 2024/131, ACM CCS 2024.} reach credentials of 80\,KB, since
improved to 54\,KB.

Set against those, our proofs are 79\,KB in identifying mode and 196\,KB across
both halves of the unlinkable one, at roughly 96 conjectured bits. That is
larger and weaker than CAPSS, and it is worth saying plainly. It is also
precisely where the theory said it would land: Jeudy et al.\footnote{ePrint
2023/560.} note in passing that anonymous credentials from STARKs are possible
but that the cost ``appears to be in the hundreds of kilobytes range.'' We
landed at 196\,KB, which is the footnote coming true rather than a surprise.

What is new is not the credential. It is where the verifier runs. Every scheme
above is verified in software, by the party being shown the credential --- a
Privacy Pass origin, a TPM verifier, a relying party. Verification is a claim
about that party's honesty. Here the verifier is a Soroban contract, so the
verdict is a fact about the ledger: the nullifier burns in consensus storage,
and the round the proof answers is derived from the ledger sequence rather than
chosen by whoever is asking. Nobody has to trust the server, because the server
is not the one deciding.

That relocation is also what explains the gap in the numbers. The 400M
instruction cap, not the cryptography, is what holds the security parameters
down; a software verifier with 20--30\,ms to spend has no such ceiling. Closing
that gap is the obvious next work, and CAPSS is where we would start.

\section{What is claimed, and what is not}

We separate these deliberately, because the value of the measurements above
depends entirely on the precision of the claims they support.

\subsection{The boundary that matters most}

Everything here concerns the \emph{proof system}. The \emph{transaction}
carrying a proof is signed with Ed25519, which is elliptic-curve and falls to
Shor. An adversary with a quantum computer does not need to break any relation
in this paper: they forge the signature on the envelope and the post-quantum
content inside becomes irrelevant. This is a post-quantum component running on
a pre-quantum chain, and the missing half is Stellar's to supply, not ours.

\subsection{Claimed}

\begin{itemize}
  \item \textbf{Post-quantum in kind}: hash-based, no elliptic curves, no
    trusted setup. Nothing in the proof system falls to Shor.
  \item \textbf{Post-quantum in strength}, at the deployed configuration:
    roughly 96 conjectured bits, near 48 against a Grover-assisted adversary.
    Both numbers rest on the standard conjectured FRI estimate rather than a
    proven bound.
  \item \textbf{Single use by consensus}: the nullifier is burned in contract
    storage, so it survives a restart of the server. It does not survive a
    redeploy of the verifier, which starts a fresh set.
  \item \textbf{Unlinkability across uses}: computational, resting on the
    hiding of the commitment under a random-oracle-style assumption with a
    248-bit blinder.
\end{itemize}

\subsection{Not claimed}

\begin{itemize}
  \item No bespoke security proof for the sponge-capacity split or the
    compression convention. The compression function is a truncated public
    permutation without feed-forward, which is \emph{provably} not collision
    resistant; no forgery follows from that in this system as far as our
    analysis goes, and both statements are recorded rather than one of them.
  \item The identifying mode does not hide which credential paid. Only the
    unlinkable mode does.
  \item The two halves of the unlinkable proof are composed by a shared public
    value, not fused into one trace, so on-chain they are two transactions.
  \item Amounts and recipients are in the clear. This hides \emph{who}, not
    \emph{how much} or \emph{to whom}.
\end{itemize}

\subsection{On priority}

An earlier draft claimed no prior x402 payment gated by an anonymous
credential. That was wrong. MicoPay\footnote{\texttt{github.com/Micopay/micopay-protocol}}
does the same shape on Stellar, independently and deployed on testnet: an agent
buys a credential over x402, spends it with a proof of membership, and a
Soroban contract burns the nullifier. We record the error here rather than
quietly narrowing the claim elsewhere.

What separates the two is the proof system. MicoPay verifies through the BN254
host functions (\texttt{g1\_msm}, \texttt{pairing\_check}); its soundness rests
on discrete log over a pairing curve, and it carries a trusted setup. Shor's
algorithm breaks that assumption. The construction here rests on a hash
function, carries no setup, and concedes only the factor of two Grover takes
out of the exponent.

The narrowed claim is therefore: no prior STARK verified directly by a Soroban
contract, and no prior \emph{post-quantum} credential gating an x402 payment.
Neither half claims the credential construction is new --- the Related work
section above places it against the literature it comes from, where it is
neither the smallest nor the strongest instance. The claim is about a verifier
that runs in consensus rather than in software.
Everything else verifying zero-knowledge proofs on Soroban today is
pairing-based --- MicoPay and StellarVeil over BN254, the SDF privacy-pool
prototype over BLS12-381, and RISC~Zero reaching Stellar only through
Nethermind's \texttt{groth16\_verifier.wasm}, a wrapper that reintroduces
pairings and a trusted setup at the last step so that no STARK ever touches the
chain.

The gap has a structural cause. CAP-0075 (Final, Protocol~25) gives Soroban a
native Poseidon2 permutation, but its \texttt{field} argument admits only
the scalar fields of BLS12-381 and BN254. Circle STARKs are cheap
precisely because they live over a 31-bit field, and no host function covers
one. A Mersenne-31 verifier must be written out in contract code, beneath the
400M-instruction cap, or it does not exist.

That is a statement about what public search surfaced on 5 August 2026, not a
proof of absence. If the remainder is wrong too, the correction belongs in the
repository and will be put there.

\section{Reproducing}

\begin{lstlisting}
git clone https://github.com/Galmanus/pq402 && cd pq402
npm install && (cd cli && npm install)
cp .env.example .env
node setup.mjs                 # generates and funds both accounts
curl -s https://channels.openzeppelin.com/testnet/gen   # a facilitator key

./demo.sh          # the identifying loop, end to end
./demo-crowd.sh    # the unlinkable one, end to end
\end{lstlisting}

One step cannot be scripted: the Circle testnet USDC faucet is captcha-gated,
and \texttt{setup.mjs} prints the address to paste there.

Running \texttt{demo-crowd.sh} twice within the same hour-long epoch produces a
refusal --- from the contract, not from the script. That refusal is the
single-use guarantee showing its face rather than being asserted.

\vspace{1.5em}
\hrule
\vspace{0.8em}
\noindent
\footnotesize
Source, transaction hashes and the full measurement log:
\url{https://github.com/Galmanus/pq402}. The Circle-STARK library the
credential relation is built from is riverrun. MIT licensed.

\end{document}
